\section{Uniform transfer of signed connection states}
\label{tr:section}
\label{sec:transfer}

The estimates of the preceding section concern annuli with four prescribed
switches at each end.  A discrete crossing test, however, can present an
arbitrarily complicated connection state to an annulus.  Moreover, the
outside of that annulus need not have a scaling limit.  We now prove
operator estimates that are uniform over these states and outside data.
The important distinction is between the two annular ends, where arbitrary
data are allowed, and an independent region between them, where the
scaling-limit arguments take place.

Throughout this section the Poisson intensity is one, the center of all
circles is fixed, and $\lambda=5/4$.  We use the localization and mesh-first
arm estimates recalled above from \cite[Sections~2--3]{V}.  The constants
implicit in choices of collars and buffers are fixed once and for all.

\subsection{Cuts, connection states, and transfer operators}

Fix a sufficiently small constant $\eta>0$.  A \emph{seam at radius $r$}
consists of the marked Poisson process in the collar
\[
             (1-\eta)r<|z|<(1+\eta)r,
\]
together with auxiliary randomness used to select a graph cycle.  A
local occupancy test, denoted $G_r$, ensures that the graph and the vertex
stars through a narrower collar are determined by these variables and
that their diameters are smaller than a fixed small fraction of $\eta r$.
We choose the test and the narrower collar with buffers at both ends, so
that
\begin{equation}
                    \Pp(G_r^c)\le C e^{-cr^2}.
\label{tr:screen}
\end{equation}
On $G_r$ choose a simple site cycle $\Gamma_r$ surrounding the center,
inside the narrower collar.  The choice uses the geometry but not the
colors.  The test and the choice are rotation equivariant; an independent
uniform angular axis may be included among the seam variables to break
ties.

For completeness, such a screen can be obtained by requiring that every
point in a closed intermediate collar be within a sufficiently small
distance, proportional to $r$, of a nucleus in the larger collar.  This
covering-distance condition is rotation invariant and has the
exponential failure bound in \eqref{tr:screen}.  Following a circle in the
site triangulation and deleting loops gives the required cycle.  Buffers
and the same covering condition determine all its vertex stars.  A
nucleus beyond either extremity of the collar cannot affect the graph
on the other side: its cell would have to pass through the screened
intermediate collar, where all cells are small.  Consequently the graph
on each side of $\Gamma_r$, including its attachments to the cycle, can
be read from the seam and the Poisson variables on that side.  Conditional
on the seam, the variables used on the two sides are independent.

The colors on $\Gamma_r$ determine an even number of switches.  An inside
connection state is a noncrossing pairing of those switches by interface
strands in the disk on the inside.  The pairing, together with the
monochromatic runs on the cycle, determines all boundary connectivity.
Closed interface loops carry no information and are not stored.  We
allow every formal noncrossing pairing as a state, whether or not it is
realized by a particular inside graph.  Such pairings respect the
alternating color convention: the endpoints of each pair have opposite
parity in cyclic order.  Gluing a pairing to the annular strands is
therefore defined without a realizability restriction.

Write $\Sigma_r$ for the seam variables and $\mathcal S_r(\Sigma_r)$ for
this finite state set on $G_r$.  A \emph{message} is a measurable field
of signed measures $v_{\Sigma_r}$ on $\mathcal S_r(\Sigma_r)$ such that
\[
 v_{\Sigma_r}(\mathcal S_r(\Sigma_r))=0,
 \qquad v_{\Sigma_r}=0\ \hbox{on }G_r^c.
\]
Its norm is
\begin{equation}
  \|v\|_r
    =\operatorname*{ess\,sup}_{\Sigma_r}
                       \|v_{\Sigma_r}\|_{\TV}.
\label{tr:message-norm}
\end{equation}
Here total variation is the sum of the absolute point masses, so that a
difference of probability measures has norm at most two.  Finite state
sets can be numbered measurably; the bounded fields of zero-mass measures,
modulo seam-law null sets, form a Banach space $\mathcal B_r$.

If the collars at radii $r<s$ are disjoint, gluing through the intervening
annulus defines a linear map
\[
                  M_{r,s}:\mathcal B_r\longrightarrow\mathcal B_s.
\]
More explicitly, conditional on $\Sigma_s$, integrate $\Sigma_r$ and
the intervening marked process, impose both seam successes, and push
$v_{\Sigma_r}$ through the annular strands.  Closed loops have weight
one.  For every fixed realization, this push is a map between finite
state spaces and contracts total variation.  Its mass remains zero.
The success restrictions are indicators in this integral, not a
conditioning followed by renormalization.  Only the output seam is
conditioned on in the definition of its message field.
Intermediate successful seams therefore give exactly the composition of
these operators, with the corresponding success indicators imposed.
This construction does not assert independence for unrestricted Voronoi
events in neighboring regions; independence is used only for the
underlying marked Poisson variables separated by the specified buffers.

Rotation and global color reversal act on seam variables and states.  A
message is \emph{isotropic} if it is fixed by the rotation action and
\emph{even} if it is fixed by color reversal.  Both symmetries are exact
for $M_{r,s}$, since the screens and cycle choices have those symmetries.

\begin{proposition}[Fixed-ratio transfer]
\label{tr:block}
There are $\delta>0$ and $C<\infty$ with the following property.  For every
sufficiently large integer $n$, set $L=2^n$.  There is a threshold $r_0(L)$
such that, for $r\ge r_0(L)$, the transfer $M=M_{r,Lr}$ satisfies
\begin{align}
 \|M\|&\le C L^{-\lambda}, \label{tr:block-size}\\
 M&=M_1+E,\qquad \operatorname{rank}M_1\le1,
       \qquad \|E\|\le C L^{-\lambda-\delta},
       \label{tr:block-rank}\\
 \|M|_{\mathcal B_r^{\mathrm{even,iso}}}\|
       &\le C L^{-2-\delta}.
       \label{tr:block-cancellation}
\end{align}
The constants $C,\delta$ do not depend on $L$ or $r$.  Every norm in the
conclusion includes the essential supremum over successful output
seams, including their colors.  The rank-one operator may depend on
$r$ and $L$.
\end{proposition}

The reduction has a simple mechanism.  Six surviving interfaces already
cost a power strictly greater than two, while four interfaces leave just
one signed connection coordinate.  We isolate an independent region in
which to obtain four-tip boundaries, verify their convergence, and then
apply the separated-end estimates to that coordinate.  Summing over the
first four-tip scales gives both the rank-one approximation and the
color--rotation cancellation.

\subsection{An independent region for the limiting argument}

Fix $L$ for the moment and rescale by $r$.  Thus the mesh tends to zero
as $r\to\infty$, whereas all radii used in the present argument lie in a
fixed finite range.  Choose auxiliary starting cycles near radii $8r$
and $Lr/8$.  Their geometry is determined in small collars separated
from both original seams.  From either starting cycle we shall explore
towards the middle of the annulus, stopping at one of finitely many
\emph{barriers}: color-independent site cycles approximating specified
concentric circles in curve order.  Their colors remain unread until
the exploration proceeds past them.  A \emph{stopping circle} means the
limiting circle of the current barrier.  All cycles may be
made chord free on the side approached by the exploration.

We impose auxiliary occupancy tests in this intervening region.  Unlike
the actual seam tests, these may require a relative mesh tending to
zero: use covering distances that tend slowly to zero in rescaled
coordinates.  For the fixed finite collection of collars under
consideration, their total failure probability is $o(1)$.  Their supports
are separated from the input and output seams.  On success, geometry
around the two starting cycles and throughout the region between them
agrees with the geometry of a full-plane Poisson sample, obtained by
adjoining independent variables outside that region.  In particular,
all these $o(1)$ failure bounds hold uniformly conditional on either or
both original seam variables.

There is also a local version of this construction.  An end experiment
stopped at a barrier uses only variables from a buffer just before its
starting cycle through a small fixed-factor enlargement past that
barrier, together with the variables between its starting cycle and its
own seam.  Its localization tests stop at that enlargement.  If two
barriers are separated by a sufficiently large fixed factor, the
supports for the two end experiments are disjoint.  Extra tests across
the middle can then be added when evaluating an actual annular
transmission.  Deleting those extra tests from an integral with bounded
weights changes it by $o(1)$, uniformly in the original seams.  It is
only after this deletion that the two end integrals are factored.
Figure~\ref{tr:figure} records the radial order and the gaps between
the supports of the two end experiments.

\begin{figure}[tb]
\centering
\begin{tikzpicture}[x=0.88cm,y=0.75cm,>=Latex,font=\small]
  \fill[black!12] (0,-0.26) rectangle (0.42,0.26);
  \fill[black!12] (11.58,-0.26) rectangle (12,0.26);
  \fill[black!5] (1.55,-0.18) rectangle (4.05,0.18);
  \fill[black!5] (7.95,-0.18) rectangle (10.45,0.18);
  \draw[thick] (0,0)--(12,0);
  \foreach \x in {0.2,2,3.7,8.3,10,11.8}
    \draw (\x,-0.14)--(\x,0.14);
  \node[align=center,below] at (0.2,-0.42) {input\\seam};
  \node[align=center,below] at (2,-0.42) {inner start\\$8r$};
  \node[align=center,below] at (3.7,-0.42) {barrier\\$a$};
  \node[align=center,below] at (8.3,-0.42) {barrier\\$b$};
  \node[align=center,below] at (10,-0.42) {outer start\\$Lr/8$};
  \node[align=center,below] at (11.8,-0.42) {output\\seam};
  \draw[->,thick] (2,0.65)--(3.7,0.65);
  \draw[->,thick] (10,0.65)--(8.3,0.65);
  \node[align=center] at (6,0.72) {unexplored annulus};
  \draw[<->] (1.2,1.5)--(10.8,1.5);
  \node[fill=white,inner sep=3pt] at (6,1.5) {independent Poisson region};
\end{tikzpicture}
\caption{The radial order of two end experiments with separated barriers,
on a schematic logarithmic scale.  Explorations start inside the transfer
block, independently of its two seams.  Each shaded end support extends
a little beyond its last barrier; the gap between those supports permits
their independent integration.  The original seam screens remain
imposed throughout.}
\label{tr:figure}
\end{figure}

The following elementary observation is the uniformity mechanism used
below.  Its purpose is to avoid any assertion of a uniform scaling
limit conditional on an atypical explored shape.

\begin{lemma}[Domination before conditioning on end shapes]
\label{tr:domination}
Suppose a contribution to $M_{r,Lr}v$ is expressed, after exploring from
the auxiliary starting cycles, as an integral of an inner scalar
coefficient, a bounded middle test, and an outer signed state measure.
If the middle test is replaced with absolute error $F\ge0$, where $F$
depends only on the independent region just described, then the error
in output norm is at most
\begin{equation}
              C\|v\|_r\,\E_{\mathrm{core}}F+o(1)\|v\|_r.
\label{tr:domination-bound}
\end{equation}
The same bound applies to discarding an event there, by taking $F$ to
be its indicator.  The expectation is under the unweighted full-plane
comparison law; the bound is uniform over successful output seams.
\end{lemma}

\begin{proof}
For a fixed realization, pushing $v_{\Sigma_r}$ to any intermediate
state set contracts total variation.  When that set has two states, its
zero-mass image has the form $\alpha(\delta_B-\delta_W)$, with
$|\alpha|\le\|v\|_r$.  Pushing the two-state difference through the other
end gives a signed measure $q$ with $\|q\|_{\TV}\le2$.  Hence the
pointwise norm of the replacement error is bounded by
$2\|v\|_rF$.  The same argument, summing absolute masses before a
two-state reduction, handles discarded events.  No regularity of the
dependence of $\alpha$ or $q$ on the seam data is needed.

On the auxiliary localization successes, $F$ is a function of marked
Poisson variables independent of both original seams and has the stated
full-plane comparison law.  Integrate the pointwise bound conditional
on the output seam.  Auxiliary failures contribute the final $o(1)$
term, also uniformly, because their supports are off the original seams.
The seam success indicators themselves are not deleted.
\end{proof}

\subsection{Exploration up to the first four-tip frontier}

First push the input message from its original seam to the inner
auxiliary starting cycle, using the marked variables in the intervening
attachment region.  The subsequent exploration uses the colored
starting cycle and the variables towards its barrier; it never consults
an attached cap state.  Dependence on the input message stays in the
pushed signed coefficients.  At the outer end the corresponding
attachment operation pushes states from its auxiliary cycle towards
the original output seam.

At the inner auxiliary cycle, reveal its colors and trace every switch
into the annulus.  A wall is followed until it pairs with another traced
wall or until its next color query would read a vertex on the current
barrier.  In the second case, stop before that query.  The stopped wall
ends at a switch edge whose next triangle has its opposite vertex on
the barrier.  No barrier colors have been read.  The analogous outer
experiment is directed inwards.  Increasing the inner barrier radius
by a factor two, or decreasing the outer barrier radius by a factor
two, continues the same wall system.  Tracing all walls in a previously
colored sample shows that the outcome does not depend on their tracing
order.

We call the unpaired stopped wall ends the \emph{surviving tips}.
Their number is even and decreases as the barrier is advanced.  The
first barrier with at most four surviving tips will be the end's
reduction barrier.  A reduction to zero or two tips kills every
zero-mass message.  A reduction to exactly four leaves precisely two
cap states.  We next verify both this assertion and the geometric
properties needed for the annular transmission limit.

\begin{lemma}[Frontier and arm counts]
\label{tr:frontier}
For either end experiment, outside an event whose probability tends to
zero as the permitted separation of tips tends to zero, mesh first,
the tips at each of finitely many fixed-scale barriers have distinct
limiting positions.  The boundary facing the unqueried barrier is a
simple site cycle whose switches are exactly these tips.  In particular:
\begin{enumerate}
\item If there are zero or two tips, the induced map on zero-mass
connection-state measures is zero.
\item If there are at least six tips, six disjoint bulk arms cross the
region from the starting cycle to the barrier, after fixed-factor
crops at the two ends.
\end{enumerate}
All probabilistic assertions are made under the unweighted exploration
law, before conditioning on a particular history.
\end{lemma}

\begin{proof}
First consider two tips approaching the same point of the barrier.
Use a graph disk patch with its base on the barrier and a half-annulus
on the exploration side.  Extend each stopped wall through its last
triangle to the base, only for this geometric argument; no barrier
color is queried.  The two wall fingers give two proper traversals of
the half-annulus, with fixed margins at its two semicircular ends.
Following the sites on their banks gives three disjoint free arms in
the sectors on the two exterior sides and between the fingers.  These
sites are on the exploration side, not on the barrier.  Short spokes
through successive wall triangles keep each bank in its sector.  The
base through the tips prevents an exterior bank from escaping around
a finger's end.  Cropping the last mesh steps gives the same count
when two stopping triangles share a barrier vertex.

This is the free-bank-shadow count for two traversals in the boundary
argument of \cite[Section~3]{V}.  The local boundary three-arm exponent
is strictly greater than one.  Covering the barrier by $O(s^{-1})$
patches of radius $s$ therefore excludes colliding tips as $s\downarrow0$,
mesh first.  The BK inequality applied to the free shadows also makes
the number of fixed-scale traversals tight.

For the frontier assertion, fill the starting side of the starting
cycle with a disk, using the sphere for the outside experiment.
Attach the queried triangles along the explored bands and add all
edges between revealed vertices.  The resulting graph is
two-connected: attaching a triangle along an existing edge preserves
two-connectivity, as does adding an edge between existing vertices.
Thus the boundary of its face containing the unqueried barrier is a
simple cycle.  The barrier and its paths towards the unexplored middle
remain in that face.  In particular each stopping triangle is incident
with this same face on its unread side.

There is a useful invariant: every revealed vertex has a revealed path
of its own color to the starting cycle.  It holds initially, and a newly
queried triangle vertex has the color of one endpoint of the incoming
heterochromatic edge, extending such a path.  Thus an added
heterochromatic edge between revealed vertices cannot belong to an
unstarted wall.  Such a wall would either close in the annulus or return
to the barrier without visiting the starting cycle; in either case it
would separate one bank from that cycle, contradicting its revealed
same-color path.  Every heterochromatic revealed-vertex edge therefore
belongs to the traced wall system.

To identify the colors along the frontier, extend the stopped walls
to the barrier as above.  These proper arcs, together with the walls
that paired back to the starting cycle, divide the explored collar
into regions with their consistent bank colors.  Every revealed vertex
is either a starting-cycle vertex or was queried in a wall triangle.
Its color agrees with the relevant bank.  An added edge cannot cross
one of these wall arcs, and the preceding invariant excludes any
untraced switch there.  Previously crossed switch edges have triangles on
both sides in the revealed bands.  The only switches exposed to the
remaining face are consequently the surviving stopping edges.

With zero or two switches there is just one noncrossing cap pairing.
Every input state gives that same pairing, so a zero-mass difference
vanishes.  With six or more surviving wall fingers, their bank shadows
give six disjoint site arms in successive sectors between the starting
cycle and the barrier.  Fixed-factor crops remove the two ends, where
different shadows could otherwise meet.  The arbitrary-color six-arm
bound is applicable, so no particular numbering or color order of the
six arms is required.
\end{proof}

The next lemma verifies the non-Jordan boundary issue rather than
presuming that scalar Cardy convergence applies to a random annulus.

\begin{lemma}[Admissibility of four-tip frontiers]
\label{tr:admissible}
At a fixed barrier, restrict the unweighted exploration law to the event
of exactly four surviving tips.  Along subsequences with good geometry
as in \cite[Section~2]{V}, the four monochromatic frontier arcs are tight
in curve order.  In every coupled subsequential limit, after constant
pauses are erased, each arc is simple and meets the stopping circle
only at its two prescribed endpoints.  The four endpoints are distinct.

For two such experiments ending at well-separated barriers, their
residual annulus satisfies the admissibility conditions of
Proposition~\ref{an:limit}.  If $t_r$ is the discrete
through-transmission probability conditional on the complete geometry
and the two end histories, and $T_r$
is the continuum transmission function evaluated on its actual
conformal data, then, for fixed $L$ and the fixed finite list of barriers,
\begin{equation}
       \E_{\mathrm{core}}\!\left[
          \boldsymbol 1_{\{\text{four tips at both ends}\}}
                              |t_r-T_r|\right]=o(1).
\label{tr:mean-limit}
\end{equation}
The expectation concerns unweighted histories, not a uniform conditional
estimate for all possible histories.
\end{lemma}

\begin{proof}
The BK and successive-excursion argument gives tightness of the four
frontier arcs.  Indeed repeated annular traversals provide disjoint
free arms; all sites used, including those on the auxiliary starting
cycle, have their original fair colors.  This is applied first at
fixed radii and then at successively smaller excursion scales.
The contact arguments below are the face-boundary arguments in
\cite[Section~3]{V}; we give the details relevant to the present
frontier.

Every frontier vertex off the starting cycle belongs to a revealed
wall band.  On its revealed side it has an opposite-colored neighbor
with a bank path back to the starting cycle.  Consequently a visit of
a frontier arc to the stopping circle other than at its designated
endpoint, with a positive excursion on both sides of the visit, gives
three free arms in a half-patch: two along the arc and one along the
opposite bank.  The same argument excludes an intermediate return to
an endpoint.  The boundary three-arm exponent and a covering of the
circle exclude all such visits.  A nonconstant interval contained in
the circle is excluded by the one-arm bound at a countable dense set
of circle points.  Constant pauses can be removed from the
parametrization.  Distinctness of the four tips follows from
Lemma~\ref{tr:frontier}.

Consider next a self-return of one monochromatic frontier arc at a
point away from the starting and stopping circles.  Separate the two
visits by a fixed positive excursion scale.  In a smaller annulus
around the proposed contact, the arc supplies four proper traversals.
Its opposite bank paths lie on the revealed side of the frontier
cycle.  They may all be joined through the starting cycle when making
the planar separation argument.  This joining set need not be
monochromatic at the starting cycle: the working annulus misses that
cycle, and within the working annulus every part of the joining set
is genuinely of the opposite color.

At least two complementary sectors have a frontier traversal on their
boundary whose revealed side faces the sector.  Each such sector
contains an opposite-colored longitudinal crossing.  To verify this assertion,
choose transverse cuts at intermediate scales with margins, leaving
both ends of that frontier traversal beyond the cuts.  If the joined
neighbor set did not cross the sector longitudinally, planar duality
would provide a transverse crosscut avoiding it.  Start this crosscut
from that frontier traversal and stop at its first subsequent hit of the
frontier cycle.  In the disk on the revealed side, the two traversal ends
and the two crosscut ends alternate along the frontier.  The crosscut
therefore separates their neighbor spokes, contradicting the joined
neighbor set, which lies on the revealed side and avoids the
crosscut.

The argument remains valid when the traversals backtrack.  Choose the
transverse cuts along intermediate graph cycles, or apply Hex duality
with such a cycle temporarily wired in the sector.  Two disjoint cuts
bound a quadrilateral whose side arcs belong to the chosen traversals
and whose interior avoids those whole proper traversals, not merely
the selected side subarcs.  Thus the crosscut's first subsequent
frontier hit lies outside the whole traversal from which it started,
so the alternating boundary order used above persists under
backtracking.  The crosscut is in the quadrilateral's interior, whereas the end spokes are beyond
the cuts.  They thus avoid the crosscut even if other portions of the
frontier enter the sector.  Cropping the margins leaves two disjoint
opposite-colored arms in different sectors, in addition to the four
arms along the monochromatic arc.  This gives six disjoint bulk arms.
The bulk exponent strictly greater than two excludes these contacts
by a covering argument.  At a self-contact on the starting circle,
the four arc traversals alone suffice: the arbitrary-color four-arm
exponent is strictly greater than one, and the possible contact
locations lie on a fixed circle.

All these exclusions are made for distinct visits with fixed positive
excursion scales before those scales decrease.  Countably many scales
and finite ball coverings make the assertions simultaneous.  They
bound joint probabilities with the four-tip event; no division by the
probability of a rare history is involved.

We now identify the remaining domain; see
Figure~\ref{fig:frontier-lobes}.  Cut along the barrier, using the
four short incidences from the tip edges through their terminal
triangles.  Between consecutive tips, the frontier arc and its
corresponding barrier interval bound a lobe polygon.  The four lobes
are attached over open intervals to the untouched region on the middle
side of the barrier.  Each lobe converges in order to a Jordan lobe,
because its frontier arc is simple and meets the barrier only at its
endpoints.  Different lobe boundaries may touch in physical projection;
their interiors remain disjoint, and their distinct accesses are
retained.  With two well-separated barriers there is a fixed open
annular region between the ends.

We verify that adjoining the lobes still gives an annulus, even when
their boundaries touch.  On either side of this open annular region,
start with the closed disk bounded by the corresponding stopping circle,
using the sphere for the outer side.  The complementary set on that
side is obtained by removing the lobe interiors and their open
attachment intervals.  Each lobe's free frontier arc remains in this
complement: it cannot enter another lobe's interior, since their Jordan
interiors are disjoint, and it meets the stopping circle only at its
own endpoints, never on another open attachment.  Connect a point of
the remaining closed set to the preserved starting-side region by a
path in the closed disk, with its interior off the stopping circle.
Whenever the path enters a lobe, replace the portion from its first
entry to its last exit by the intervening subarc of that lobe's free
frontier.  The entry and exit lie on this free arc.  The replacement
avoids every lobe interior and every open attachment, so it cannot
undo previous replacements.  There are only four lobes; after finitely
many replacements the path is wholly in the complementary set.
Thus each side has a connected complement.  The two such nondegenerate
continua are separated by the open annular region and exhaust the
complement of the residual domain.  The domain is therefore doubly
connected, with its two boundary components physically apart.

Lemma~\ref{an:back-lobes} now gives uniform convergence
of the closed-annulus parametrizations with their prime-end order.
The boundary traversal consists of finitely many individually simple
arcs.  These are exactly the hypotheses required by
Proposition~\ref{an:limit}.

Finally extract, from any sequence of scales, a further sequence on
which the localized full-plane geometries are good in the sense of
\cite[Section~2]{V}.  Conditional on those geometries, couple convergent
subsequences of the frontier histories just described.  All unqueried
colors in the residual face remain independent and fair.  Apply
Proposition~\ref{an:limit} path by path to the converging graph annuli.
The continuum function on the actual approximating conformal data
also converges, by its continuity.  Bounded convergence proves
\eqref{tr:mean-limit} along the further sequence and hence along the
original sequence.  Unused colors on the removed sides of the
frontiers may be revealed or integrated in their own end variables;
they do not affect this argument.
\end{proof}

\begin{figure}[tb]
\centering
\begin{tikzpicture}[font=\footnotesize,>=Latex,line cap=round,
  bank/.style={very thick},
  barrier/.style={densely dashed,gray!80},
  fluid/.style={fill=blue!8}]
  \begin{scope}[xshift=0.3cm]
    \fill[blue!8] (0,0) circle (2.65);
    \fill[black!9]
      (2,0) .. controls (1.25,.25) and (.25,1.25) .. (0,2)
      .. controls (-.25,1.25) and (-1.25,.25) .. (-2,0)
      .. controls (-1.25,-.25) and (-.25,-1.25) .. (0,-2)
      .. controls (.25,-1.25) and (1.25,-.25) .. cycle;
    \draw[barrier] (0,0) circle (2);
    \draw[bank]
      (2,0) .. controls (1.25,.25) and (.25,1.25) .. (0,2);
    \draw[bank,gray!55]
      (0,2) .. controls (-.25,1.25) and (-1.25,.25) .. (-2,0);
    \draw[bank]
      (-2,0) .. controls (-1.25,-.25) and (-.25,-1.25) .. (0,-2);
    \draw[bank,gray!55]
      (0,-2) .. controls (.25,-1.25) and (1.25,-.25) .. (2,0);
    \foreach \p in {(2,0),(0,2),(-2,0),(0,-2)}
      \fill \p circle (2pt);
    \draw[gray!75] (0,0) circle (.58);
    \node[align=center,font=\scriptsize] at (0,0) {starting\\side};
    \node at (1.15,1.15) {$L_1$};
    \node at (-1.15,1.15) {$L_2$};
    \node at (-1.15,-1.15) {$L_3$};
    \node at (1.15,-1.15) {$L_4$};
    \node[align=center] at (0,2.97) {unqueried annular back};
    \draw[->,thin] (.78,2.87)--(1.15,2.18);
    \node[align=left,anchor=west] at (2.8,.95)
      {four-tip\\frontier};
    \draw[->,thin] (2.78,.77)--(1.43,.30);
    \node[align=left,anchor=west] at (2.8,-.80)
      {open attachment\\interval};
    \draw[->,thin] (2.78,-.70)--(1.74,-.97);
    \node[align=center] at (0,-3.05)
      {Dashed circle: stopping barrier.\\
       Its open intervals lie inside the glued domain.};
  \end{scope}

  \begin{scope}[xshift=8.85cm,yshift=.20cm]
    \node[align=center] at (0,2.28) {Two distinct banks at a contact};
    \fill[blue!8]
      (-1.55,1.70)--(1.55,1.70)--(1.55,.85)
      .. controls (.85,.85) and (.50,0) .. (0,0)
      .. controls (-.50,0) and (-.85,.85) .. (-1.55,.85)--cycle;
    \fill[blue!8]
      (-1.55,-1.70)--(1.55,-1.70)--(1.55,-.85)
      .. controls (.85,-.85) and (.50,0) .. (0,0)
      .. controls (-.50,0) and (-.85,-.85) .. (-1.55,-.85)--cycle;
    \draw[bank]
      (-1.55,.85) .. controls (-.85,.85) and (-.50,0) .. (0,0)
      .. controls (.50,0) and (.85,.85) .. (1.55,.85);
    \draw[bank,gray!55]
      (-1.55,-.85) .. controls (-.85,-.85) and (-.50,0) .. (0,0)
      .. controls (.50,0) and (.85,-.85) .. (1.55,-.85);
    \fill (0,0) circle (2.2pt);
    \node at (0,1.30) {first access};
    \node at (0,-1.30) {second access};
    \draw[->,thin] (0,.96)--(0,.20);
    \draw[->,thin] (0,-.96)--(0,-.20);
    \node[align=center] at (0,-2.30)
      {The contact is boundary,\\not a gluing interval.};
  \end{scope}
\end{tikzpicture}
\caption{The decomposition at an inner four-tip frontier. The four
monochromatic frontier arcs, drawn alternately black and gray, bound
Jordan lobes $L_1,\ldots,L_4$ with the stopping barrier. Each lobe is
attached to the untouched annular back along an open barrier interval;
the common shading denotes the resulting unqueried domain. The starting
side remains in the complement. At the other end the construction is
reversed. The inset records two accesses at a possible contact between
different banks: the contact is not added to the fluid and does not join
the accesses. All curves are schematic. Actual banks may be rough and
may touch each other, although each individual frontier arc is simple.}
\label{fig:frontier-lobes}
\end{figure}


The probabilistic conclusions of the last two lemmas all concern the
independent region.  Lemma~\ref{tr:domination} therefore makes their
discarded probabilities and the error in \eqref{tr:mean-limit} uniform
for the operator norm, even when a message depends arbitrarily on the
seam and on its attachment information.  This is the reason the
explorations start at auxiliary cycles rather than at the original
seams.

\subsection{The first-reduction decomposition}

Index the inner barriers by radii
\[
             a_j=8\cdot2^j r,
\qquad j=1,2,\ldots,
\]
and the outer barriers by
\[
             b_k=2^{-k}Lr/8,
\qquad k=1,2,\ldots.
\]
Each list is finite, ending before the opposite starting region.  Fix
a large integer $d$ so that
\begin{equation}
                  j+k\le n-d
\label{tr:used-pairs}
\end{equation}
ensures disjoint buffered end supports and sufficient separation for
Proposition~\ref{sh:separated-ends}.  A \emph{used pair} is a pair
satisfying \eqref{tr:used-pairs} at which the two experiments first
reduce to exactly four tips.  A first reduction to fewer than four
contributes zero by Lemma~\ref{tr:frontier}.

Choose $\tau>0$ smaller than the excess of the mesh-first six-arm
exponent over two.  Let $\rho_j^{\rm in}$ and $\rho_k^{\rm out}$ be the
subprobability laws of the corresponding first four-tip histories,
including local screens and the variables attaching each history to its
original seam.  The inner law integrates the input seam.  All outer
laws and outer integrals below are conditional on the fixed output seam
$\Sigma_{Lr}$; in particular they define fields over that seam, not
averages of it.  The dependence on seam variables in the
attachments will be kept in the scalar and signed measure factors
below, rather than in any assertion about typical shapes.  We first
bound the masses of these unweighted history laws; pointwise bounded
message factors will then preserve the bounds up to constants.

\begin{lemma}[First-reduction tails]
\label{tr:tails}
For fixed $L$, the end-history subprobability masses satisfy, in the
mesh-first limit,
\begin{equation}
 \|\rho_j^{\rm in}\|_{\TV}
 \operatorname*{ess\,sup}_{\Sigma_{Lr}}
         \|\rho_k^{\rm out}(\cdot\mid\Sigma_{Lr})\|_{\TV}
      \le C 2^{-(2+\tau)j}2^{-(2+\tau)k}+o(1)
\label{tr:history-weight}
\end{equation}
for each used pair $(j,k)$.  Consequently, after inserting arbitrary
middle tests of modulus at most one, the same bound up to a constant holds
for the corresponding operator contribution on inputs of norm at most
one.  The part of the transfer not represented
by used pairs has norm at most
\begin{equation}
             C(n+1)^2L^{-2-\tau}+o(1).
\label{tr:unused}
\end{equation}
Both estimates are uniform for input messages of norm at most one and
for successful output seams.
\end{lemma}

\begin{proof}
Before the first reduction at level $j$, there are at least six tips
at the preceding level.  Lemma~\ref{tr:frontier} then requires six arms
from the starting region to that level.  The mesh-first six-arm bound
gives $C2^{-(2+\tau)j}$, with a changed constant for the first few
levels.  The same statement holds at the outer end, uniformly in its
original seam: the necessary six-arm event is read in the independent
region, while attachment variables and success indicators have total
conditional mass at most one.  Multiplying the two separate mass
bounds proves \eqref{tr:history-weight}.  For used pairs the physical
end experiments also have disjoint buffered supports, so their joint
law is the product used in the operator contribution.  Arbitrary input
weights are bounded
pointwise by their total variation, as in
Lemma~\ref{tr:domination}; the number of switches on the input seam
does not enter the estimate.

To bound the remainder, run the inner exploration to its longest
allowed level.  If it never reduces, it has paid a six-arm probability
over $n-O(1)$ dyadic levels.  Otherwise let $j$ be its first reduction
level.  If the reduction has fewer than four tips, its contribution
is zero.  In the remaining case run the outer experiment as far as
\eqref{tr:used-pairs} allows, possibly through no additional level.
If it does not reduce in time, the two experiments together have
paid six-arm probabilities over at least $n-O(1)$ levels.  Their
buffered supports remain disjoint.  Summing the resulting bounds over
the at most $O((n+1)^2)$ possible terminal levels proves
\eqref{tr:unused}.  The finitely many auxiliary localization failures
and near-coincidences of tips cost $o(1)$ for fixed $L$, uniformly by
Lemma~\ref{tr:domination}.
\end{proof}

At a four-tip inner frontier let $B$ and $W$ denote respectively the
black-join and white-join cap states.  The pushed input message is
\begin{equation}
                 \alpha_j(v;H_i)(\delta_B-\delta_W),
                 \qquad |\alpha_j(v;H_i)|\le\|v\|_r,
\label{tr:inner-coefficient}
\end{equation}
where $H_i$ includes the inner history and its attachment to the input
seam.  The coefficient is linear in $v$.  Similarly, pushing
$\delta_B-\delta_W$ from an outer four-tip frontier to the output seam
gives a signed measure $q_k(H_o)$ satisfying
\begin{equation}
                         \|q_k(H_o)\|_{\TV}\le2.
\label{tr:outer-coefficient}
\end{equation}
Lemma~\ref{an:caps} says that the residual annulus sends the first
cap difference to the second multiplied by its through-transmission
probability.  Thus the used-pair contribution has the form
\begin{equation}
 \int \alpha_j(v;H_i)\,
       \E_{\rm mid}\!\left[
          \boldsymbol1_{G_{\rm aux}}t_r\mid H_i,H_o\right]
       q_k(H_o)\,d\rho_j^{\rm in}(H_i)d\rho_k^{\rm out}(H_o),
\label{tr:integral}
\end{equation}
where $G_{\rm aux}$ denotes the additional middle screens and
$\E_{\rm mid}$ integrates the middle Poisson geometry conditional on
the two end histories.  The probability $t_r$ already integrates the
fresh colors in the remaining face, conditional on all geometry.
The end laws are a product because their buffered supports are
disjoint.  The conditional factor in square brackets still couples
the ends, and must not yet be factored.

On the auxiliary localization successes, the residual frontiers and
both $t_r$ and $T_r$ depend only on the auxiliary-cycle explorations and
the intervening independent geometry.  Attachments to the original
seams enter only the signed coefficients.  The error $|t_r-T_r|$ thus
has exactly the independent-region measurability required by
Lemma~\ref{tr:domination}.

By Lemmas~\ref{tr:admissible} and~\ref{tr:domination}, replace $t_r$ by
$T_r$ in \eqref{tr:integral} with $o(1)\|v\|_r$ error in output norm.
For fixed $n$, this error may be summed over all used pairs.  The
frontiers themselves are still the actual graph cycles.  Their
single-end conformal maps and capacities can therefore be used
directly in Proposition~\ref{sh:separated-ends}, without asserting
smoothness at contacts of a limiting boundary.

\subsection{Rank-one factorization and signed cancellation}

We now complete Proposition~\ref{tr:block}.  An inner frontier at
level $j$ surrounds the center, stays inside radius comparable to
$a_j$, and reaches that radius at its tips.  Its logarithmic capacity
$l_i$ is therefore comparable to $a_j$.  An outer frontier at level
$k$ surrounds a disk of radius comparable to $b_k$ and has a point
at distance comparable to $b_k$; its conformal radius $l_o$ is
comparable to $b_k$.  Write $y^0$ and $x^0$ for their single-end
switch angles.  By \eqref{sh:rank-one},
\begin{equation}
 T_r(H_i,H_o)=K\,U(H_i)V(H_o)
                       +O\big((a_j/b_k)^{\lambda+\sigma_1}\big),
\label{tr:separated-formula}
\end{equation}
where
\[
 U(H_i)=l_i^\lambda J(y^0)^{1/3},
 \qquad
 V(H_o)=l_o^{-\lambda}J(x^0)^{1/3}.
\]
These factors depend on their own end histories only, and
\begin{equation}
             |U(H_i)V(H_o)|
                  \le C L^{-\lambda}2^{\lambda(j+k)}.
\label{tr:factor-size}
\end{equation}

For the separated leading term, delete the extra middle screens.
This changes the integral by $o(1)\|v\|_r$ for fixed $L$: the screens
use independent interior variables and fail with probability $o(1)$,
while all weights at these fixed scales are bounded.  Keep the original
seam screens and the local end screens.  The remaining end experiments
now use disjoint Poisson supports for every used pair.  Define the
scalar functionals and output messages
\[
 u_j(v)=\int \alpha_j(v;H_i)U(H_i)\,d\rho_j^{\rm in}(H_i),
 \qquad
 v_k=\int V(H_o)q_k(H_o)\,d\rho_k^{\rm out}(H_o).
\]
The leading contribution of that pair is exactly $K u_j(v)v_k$.
The same definitions make sense separately for every level in each
finite list, whether or not a pair of levels has disjoint supports.
For such additional pairs we use the product of the two separate
integrals, not a claim about simultaneous experiments.

Set
\begin{equation}
       M_1v=K\left(\sum_j u_j(v)\right)\left(\sum_k v_k\right).
\label{tr:rank-one-definition}
\end{equation}
This operator has rank at most one.  By the separate six-arm bounds
and \eqref{tr:factor-size}, its norm is at most
\begin{equation}
 C L^{-\lambda}
       \sum_{j,k}2^{-(2+\tau-\lambda)(j+k)}
                   \le C L^{-\lambda}.
\label{tr:rank-one-size}
\end{equation}
The extra product terms introduced by allowing $j+k>n-d$ have norm
at most
\begin{equation}
 C L^{-\lambda}
      \sum_{j+k>n-d}2^{-(2+\tau-\lambda)(j+k)}
                   \le C(n+1)^2L^{-2-\tau}.
\label{tr:artificial-pairs}
\end{equation}
This numerical estimate requires no independence for overlapping
actual experiments.

Decrease $\sigma_1$ if necessary so that
$\lambda+\sigma_1<2+\tau$.  The errors in
\eqref{tr:separated-formula}, summed using
Lemma~\ref{tr:tails}, have norm at most
\begin{equation}
 C L^{-\lambda-\sigma_1}
      \sum_{j+k\le n-d}
             2^{-(2+\tau-\lambda-\sigma_1)(j+k)}+o(1)
       \le C L^{-\lambda-\sigma_1}+o(1).
\label{tr:rank-error-sum}
\end{equation}
Together with \eqref{tr:unused} and
\eqref{tr:artificial-pairs}, this proves the first two estimates of
Proposition~\ref{tr:block}, with a smaller positive exponent saving,
in the mesh-first limit.  All errors are measured in output total
variation, not merely after testing a fixed continuum event.

It remains to prove the stronger estimate for an even isotropic input
$v$.  After the continuum replacement, delete the extra middle screens
as above.  Only $T_r$ then depends on both ends; all other factors
belong to their separate, independently integrated end laws.
Rotate all inner end data, including the seam and the revealed
history.  By equivariance of $v$, the coefficient
$\alpha_j(v;H_i)$ in \eqref{tr:inner-coefficient} is unchanged.  Reverse
all colors of the inner data instead.  The message is still unchanged
as a geometrically transformed signed measure, but the reference
vector $\delta_B-\delta_W$ changes sign.  Thus
\begin{equation}
 \alpha_j(v;\operatorname{rot}_\theta H_i)=\alpha_j(v;H_i),
 \qquad
 \alpha_j(v;\operatorname{flip} H_i)=-\alpha_j(v;H_i).
\label{tr:coefficient-symmetry}
\end{equation}
The inner end law is rotation and color-reversal invariant.  The outer
measure factor is unaffected by either transformation of the inner
data.  Labels for the four tips may be chosen independently of colors
using the auxiliary angular axis, and transported under rotations.
For each fixed outer history $H_o$, color reversal gives the identity
\begin{equation}
\begin{aligned}
 &\int \alpha_j(v;H_i)T_r(H_i,H_o)\,d\rho_j^{\rm in}(H_i)\\
 &\quad=\frac12\int \alpha_j(v;H_i)
    \big[T_r(H_i,H_o)-T_r(\operatorname{flip}H_i,H_o)\big]
                                   \,d\rho_j^{\rm in}(H_i).
\end{aligned}
\label{tr:flip-integral}
\end{equation}
Now average rotations of the inner data in the right side.  Since
$\alpha_j$ and the inner law are rotation invariant, its contribution
is one half of the integral against the rotation average of
$T^+-T^-$, with its appropriate inner prescriptions.  In particular
one must retain the sign; replacing this difference by an unsigned
four-arm event would lose the required exponent.

Equation~\eqref{sh:rotation-average} bounds the averaged difference
by $C(a_j/b_k)^{2+\sigma_2}$.  The local end experiments and screens
can be chosen rotation and color-reversal invariant.  Decrease
$\sigma_2$ so that $0<\sigma_2<\tau$.  The used-pair contributions
are bounded in total variation by
\begin{equation}
 C L^{-2-\sigma_2}
    \sum_{j+k\le n-d}2^{-(\tau-\sigma_2)(j+k)}+o(1)
                  \le C L^{-2-\sigma_2}+o(1).
\label{tr:signed-sum}
\end{equation}
Adding the unused contribution \eqref{tr:unused} proves the desired
power strictly greater than two.

Finally choose $\delta>0$ smaller than all the positive savings just
obtained, decreasing it to absorb the polynomial factors in $n$.
For each fixed sufficiently large $n$, choose $r_0(2^n)$ so that every
$o(1)$ term is smaller than the required fixed power of $L$.  This
proves \eqref{tr:block-size}--\eqref{tr:block-cancellation}.
Notice the order: the ratio is chosen first and the scale threshold
afterwards.  No rate of convergence in Cardy's formula, in the
good-sequence extraction, or near a nonsimple limiting bank is used.

\subsection{Iteration and failed screens}

The block estimates have two consequences.  The first uses the signed
symmetry estimate to control a Poisson insertion.  The second uses the
rank-one approximation to compare two responses at different distant
boundaries.

\begin{corollary}[Fair-insertion message]
\label{tr:insertion}
There are $0<\nu<1$ and $C<\infty$ such that the following holds.  Insert
a nucleus at the center of a cut, with an independent fair color.  At
a seam of radius $s$, take the difference of the inside state laws
``with this insertion'' minus ``without this insertion,'' conditional
on the seam and with its screen success imposed.  Its message
$I_s$ satisfies
\begin{equation}
                            \|I_s\|_s\le C s^{-2-\nu}
\label{tr:insertion-bound}
\end{equation}
for all sufficiently large $s$.  The estimate remains valid in the
presence of fixed data outside the screened collar that do not enter
its inside.  The seam tests and the selected cycle are those of the
base process, without the insertion.
\end{corollary}

\begin{proof}
Choose a ratio $L$ from Proposition~\ref{tr:block} so large that the
constant per block in \eqref{tr:block-cancellation} is absorbed by a
smaller power saving.  Choose $s_*>r_0(L)$ sufficiently large.  For
given $s$, use radii
\[
             r_j=L^j r_0,\qquad 0\le j\le N,
             \qquad r_N=s,\quad s_*\le r_0<Ls_*.
\]
All screens use the base process.  At a successful seam, the inserted
point cannot reach or alter its cycle, by localization.  The two
inside state laws consequently have the same state space and each
has mass one.  Their difference has zero mass, norm at most two,
and is even and isotropic.

If every seam is successful, iterating
\eqref{tr:block-cancellation} bounds the message by
$C(s/r_0)^{-2-\nu}$ for some $0<\nu<1$.  Since $r_0$ stays in a fixed
compact range of positive scales, this gives
\eqref{tr:insertion-bound} for that contribution.

For the remaining contribution, group configurations by the largest
failed radius $v=r_j<s$.  The exact partition, with the final success
kept in every term, is
\[
 \boldsymbol1_{G_{r_N}}
   =\prod_{i=0}^{N}\boldsymbol1_{G_{r_i}}
    +\sum_{j=0}^{N-1}\boldsymbol1_{G_{r_j}^c}
                      \prod_{i=j+1}^{N}\boldsymbol1_{G_{r_i}}.
\]
Form a fresh difference of inside state
laws at the next cut $Lv=r_{j+1}$, multiplying both laws by the same
base-process indicator $\boldsymbol1_{G_v^c}$.  Earlier screen
outcomes are left unrestricted.  This weighted difference still has
zero mass for each successful seam at $Lv$: for each common base
configuration the two state laws have equal total mass.  It is also
even and isotropic.  The collar at $v$ is disjoint from that at $Lv$,
so conditional on the latter seam its norm is at most
\begin{equation}
                2\Pp(G_v^c\mid\Sigma_{Lv})
                           \le C e^{-cv^2}.
\label{tr:reset}
\end{equation}
Success at $Lv$ ensures that both configurations use the same cut
there, independently of what happened at smaller radii.

Transfer this reset message through the successful cuts above $Lv$.
Those success restrictions ensure that $v$ is indeed the largest
failed radius; no restriction on earlier failures was needed.
The sum of these contributions is bounded by
\[
 C\sum_{v\in\{r_0,\ldots,r_{N-1}\}}
             e^{-cv^2}\left(\frac{s}{Lv}\right)^{-2-\nu}
       \le C s^{-2-\nu}.
\]
The final inequality follows by summing the exponentially decaying
weights times $v^{2+\nu}$ over the geometric sequence.  This proves
the claim.  Outside data beyond the final screened collar have no
effect on the inside law or on the argument, which proves the stated
stability.
\end{proof}

For a message space $\mathcal B_t$, the projective norm on its algebraic
tensor square is
\[
 \|z\|_\pi
   =\inf\left\{\sum_{i=1}^m\|u_i\|_t\|v_i\|_t:
          z=\sum_{i=1}^m u_i\otimes v_i,\quad m<\infty\right\}.
\]
We may complete in this norm.  Bounded linear functionals $\ell,\ell'$
on $\mathcal B_t$ satisfy
\[
         |(\ell\otimes\ell')(z)|
                \le\|\ell\|\,\|\ell'\|\,\|z\|_\pi.
\]
In particular, evaluating $u\otimes v-v\otimes u$ gives the determinant
$\ell(u)\ell'(v)-\ell(v)\ell'(u)$.

\begin{corollary}[Contraction of antisymmetric products]
\label{tr:wedge}
There are $\nu'>0$, $s_0<\infty$, and a fixed ratio $L>1$ such that,
for $r\ge s_0$ and $s/r=L^N$, the transfer
\[
          \mathcal M=M_{s/L,s}\cdots M_{r,Lr}
\]
with all intermediate seam successes imposed satisfies
\begin{equation}
 \|\mathcal M u\wedge\mathcal M v\|_\pi
       \le C(s/r)^{-2\lambda-\nu'}\|u\|_r\|v\|_r
\label{tr:wedge-bound}
\end{equation}
for any two zero-mass input messages $u,v\in\mathcal B_r$.  Here
$u\wedge v=u\otimes v-v\otimes u$ and $\|\cdot\|_\pi$ is the
projective tensor norm.  Fixed-factor unused intervals at the two
ends change only the constant.  No assertion of this contraction
through a failed intermediate cut is made.
\end{corollary}

\begin{proof}
For one block write $M=M_1+E$ as in
\eqref{tr:block-rank}.  Since $M_1$ has rank at most one,
$M_1\otimes M_1$ vanishes on antisymmetric tensors.  On that subspace,
\[
 \|M\otimes M\|
        \le 2\|M_1\|\,\|E\|+\|E\|^2
        \le C L^{-2\lambda-\delta}.
\]
Choose $L$ large enough to absorb $C$ into a smaller exponent saving
$\nu'>0$, simultaneously with its choice in
Corollary~\ref{tr:insertion}.  Iterating the estimate and using
$\|u\wedge v\|_\pi\le2\|u\|_r\|v\|_r$ gives
\eqref{tr:wedge-bound}.  The argument holds first for finite tensor
sums and then for their projective-norm completion.  Ordinary
Markov pushes contract total variation, so fixed-factor omitted end
intervals affect only the constant after rounding to the ratio-$L$
grid.
\end{proof}
